% Modified from supplied SCU course-paper template. Original licence preserved in ../LICENSE.
% Changes: English body labels, course title, long-name cover fitting.
% !TEX program = xelatex
% ============================================================
%  四川大学 · 研究生课程论文 LaTeX 模板
%  ----------------------------------------------------------
%  编译方式：XeLaTeX（Overleaf: Menu → Compiler 选 XeLaTeX；
%            本模板已附带 latexmkrc，Overleaf 通常会自动识别）
%  含参考文献的完整编译顺序：
%      xelatex → bibtex → xelatex → xelatex
%  排版规范改编自四川大学研究生学位论文模板 scuthesis2020
%  （https://github.com/kevinleeex/scu_thesis_2020，GPL 协议）
% ============================================================
\documentclass[a4paper,UTF8,zihao=-4,fontset=fandol]{ctexart}
% fontset=fandol：Overleaf / 任意系统免装中文字体即可编译；
% 本地 Windows 上如想用微软雅黑/宋体，可改为 fontset=windows。

% ---- 编译器保险：本模板只能用 XeLaTeX，误选 pdfLaTeX 时给出中文提示 ----
\usepackage{iftex}
\ifPDFTeX
  \PackageError{scu-course-paper}
    {本中文模板必须使用 XeLaTeX 编译，当前检测到的是 pdfLaTeX！\MessageBreak
     Overleaf：Menu -> Compiler -> 选择 XeLaTeX；\MessageBreak
     本地编译：使用命令 xelatex main.tex（不要用 pdflatex）}
    {切换编译器后重新编译即可。}
\fi

% ----------------- 封面 logo 开关 -----------------
\newif\ifcolorlogo
\colorlogotrue          % 彩色川大校徽；需要黑白版改为 \colorlogofalse

% ===================== 宏包 =====================
\usepackage{geometry}
\geometry{left=2.5cm,right=2.5cm,top=2.75cm,bottom=3cm,headsep=0.6cm,footskip=1cm}

\usepackage{graphicx}
\usepackage{caption}
\usepackage{subcaption}
\captionsetup{font=small,labelfont=bf,labelsep=quad}
\captionsetup[figure]{name=Figure}
\captionsetup[table]{name=Table}

\usepackage{amsmath,amssymb,amsfonts,amsthm,mathtools}
\usepackage{booktabs,longtable,multirow,array}
\usepackage{enumitem}
\setlist[itemize]{itemsep=2pt,parsep=0pt,topsep=4pt}
\setlist[enumerate]{itemsep=2pt,parsep=0pt,topsep=4pt}
\usepackage{float}

% 算法（中文标签）

\usepackage{xcolor}
% Unused TikZ illustration dependencies removed.

\usepackage{fancyhdr}
\usepackage[sort&compress]{gbt7714}   % GB/T 7714-2015 参考文献
\usepackage{hyperref}
\hypersetup{hidelinks}                % 打印友好，不显示链接彩色框；
                                      % 想要彩色链接可改为 colorlinks=true
\usepackage{cleveref}                 % 智能引用：\cref{fig:x} → 图 3-1

% 封面下划线（与学位论文模板同款；xeCJKfntef 仅 XeLaTeX 可用）
\ifXeTeX\usepackage{xeCJKfntef}\fi

% ===================== 基础格式 =====================
\linespread{1.5}                       % 1.5 倍行距
\setlength{\parindent}{2em}           % 首行缩进两字符
\setlength{\parskip}{2pt}
\setcounter{secnumdepth}{3}
\setcounter{tocdepth}{3}

% 图表、公式、定理、算法按“节”编号：图 1-1、式 (2-1)
\numberwithin{equation}{section}
\numberwithin{figure}{section}
\numberwithin{table}{section}
\renewcommand{\theequation}{\arabic{section}-\arabic{equation}}
\renewcommand{\thefigure}{\arabic{section}-\arabic{figure}}
\renewcommand{\thetable}{\arabic{section}-\arabic{table}}

% 定理类环境
\theoremstyle{plain}
\newtheorem{theorem}{定理}[section]
\newtheorem{lemma}[theorem]{引理}
\newtheorem{proposition}[theorem]{命题}
\newtheorem{corollary}[theorem]{推论}
\theoremstyle{definition}
\newtheorem{definition}[theorem]{定义}
\newtheorem{example}[theorem]{例}
\theoremstyle{remark}
\newtheorem{remark}[theorem]{注}
\renewcommand{\proofname}{证明}
\renewcommand{\thetheorem}{\arabic{section}-\arabic{theorem}}

% 各级标题：节（三号黑体居中）/ 小节（四号黑体）/ 小小节（小四黑体）
\ctexset{
  section/aftername={\hspace{\ccwd}},
  section={format=\centering\zihao{3}\heiti\bfseries,beforeskip=18pt,afterskip=18pt},
  subsection/aftername={\hspace{\ccwd}},
  subsection={format=\raggedright\zihao{4}\heiti\bfseries,beforeskip=18pt,afterskip=6pt},
  subsubsection/aftername={\hspace{\ccwd}},
  subsubsection={format=\raggedright\zihao{-4}\heiti\bfseries,beforeskip=12pt,afterskip=6pt},
  paragraph={format=\raggedright\zihao{-4}\kaishu,beforeskip=6pt,afterskip=0pt}
}
\renewcommand{\contentsname}{Contents}

% 页眉：《课程名称》课程论文；页脚居中页码
\makeatletter
\fancypagestyle{scuhead}{%
  \fancyhf{}%
  \fancyhead[C]{\zihao{-5}\songti 《\@coursename》课程论文}%
  \fancyfoot[C]{\zihao{-5}\songti\thepage}%
  \renewcommand{\headrulewidth}{0.8pt}}
\fancypagestyle{plain}{%
  \fancyhf{}%
  \fancyfoot[C]{\zihao{-5}\songti\thepage}%
  \renewcommand{\headrulewidth}{0pt}}
\makeatother
\pagestyle{scuhead}

% cleveref 中文引用名
\crefname{section}{节}{节}
\Crefname{section}{第}{节}
\crefname{figure}{图}{图}
\Crefname{figure}{图}{图}
\crefname{table}{表}{表}
\Crefname{table}{表}{表}
\crefname{equation}{式}{式}
\Crefname{equation}{式}{式}
\crefname{algorithm}{算法}{算法}
\Crefname{algorithm}{算法}{算法}
\creflabelformat{equation}{(#2#1#3)}

% ===================== 封面信息接口 =====================
\makeatletter
\newcommand{\papertype}[1]{\def\@papertype{#1}}
\newcommand{\coursename}[1]{\def\@coursename{#1}}
\newcommand{\papertitle}[1]{\def\@papertitle{#1}}
\newcommand{\college}[1]{\def\@college{#1}}
\newcommand{\majorname}[1]{\def\@majorname{#1}}
\newcommand{\studentname}[1]{\def\@studentname{#1}}
\newcommand{\studentid}[1]{\def\@studentid{#1}}
\newcommand{\teachername}[1]{\def\@teachername{#1}}
\newcommand{\termpaper}[1]{\def\@term{#1}}
\def\@papertype{研究生课程论文}
\def\@coursename{}
\def\@papertitle{}
\def\@college{}
\def\@majorname{}
\def\@studentname{}
\def\@studentid{}
\def\@teachername{}
\def\@term{}

% 封面填空下划线（与学位论文模板一致）
\newcommand{\fillinblank}[2]{%
  \CJKunderline[thickness=0.8pt,depth=0.45em]{\makebox[#1]{#2}}}

\newcommand{\makecover}{%
  \begin{titlepage}
    \thispagestyle{empty}
    \begin{center}
      \vspace*{-0.3cm}
      \ifcolorlogo
        \includegraphics[width=9.2cm]{figures/SCU_TITLE.png}%
      \else
        \includegraphics[width=9.2cm]{figures/SCU_TITLE_BW.png}%
      \fi
      \\[1.5cm]
      {\zihao{1}\heiti\bfseries \@papertype}
      \\[2.6cm]
      {\zihao{-3}
       \renewcommand{\arraystretch}{1.4}
       \begin{tabular}{@{}r@{\hspace{0.5em}}l@{}}
         {\heiti\makebox[2.8cm][s]{论文题目：}} & \fillinblank{7.4cm}{\songti\@papertitle}\\
         {\heiti\makebox[2.8cm][s]{课程名称：}} & \fillinblank{7.4cm}{\songti\@coursename}\\
         {\heiti\makebox[2.8cm][s]{学　　院：}} & \fillinblank{7.4cm}{\songti\@college}\\
         {\heiti\makebox[2.8cm][s]{专业名称：}} & \fillinblank{7.4cm}{\songti\@majorname}\\
         {\heiti\makebox[2.8cm][s]{学生姓名：}} & \fillinblank{7.4cm}{\songti\@studentname}\\
         {\heiti\makebox[2.8cm][s]{学　　号：}} & \fillinblank{7.4cm}{\songti\@studentid}\\
         {\heiti\makebox[2.8cm][s]{任课教师：}} & \fillinblank{7.4cm}{\songti\@teachername}\\
       \end{tabular}}
      \vfill
      {\zihao{4}\songti \@term}
      \vspace{1.0cm}
    \end{center}
  \end{titlepage}}
\makeatother


\setlength{\headheight}{16pt}
\emergencystretch=2em
